% chktex-file 0
\documentclass{ctexart}
\usepackage{graphicx} % Required for inserting images
\usepackage[margin=2.5cm]{geometry}
\usepackage{booktabs}
\usepackage{float}
\usepackage{soulutf8}
\usepackage{enumitem}
\usepackage{pifont}
\usepackage{pgfplots}
\usepackage{longtable}
\pgfplotsset{compat=newest}
\usepackage{tikz}
\usepackage{xcolor} % 用于颜色设置
\usepackage{listings} % 用于代码块
\usepackage{fontspec} % 关键：用于使用系统字体 (XeLaTeX/LuaLaTeX)
% Linux 环境中没有 Windows 宋体时，使用可用的中文衬线字体作为宋体替代
\setCJKmainfont{Noto Serif CJK SC}
\usepackage{authblk}
\usepackage[
    colorlinks=true,      % 使用颜色而非方框
    linkcolor=blue,       % 内部链接颜色
    citecolor=green,      % 引用链接颜色
    urlcolor=red,         % URL 链接颜色
    bookmarks=true,       % 生成 PDF 书签
    bookmarksopen=true,   % 默认展开书签
    pdfstartview=FitH     % 页面宽度适应窗口
]{hyperref}

\sethlcolor{yellow}%设置所需高亮颜色
\usepackage[most]{tcolorbox}
\tcbuselibrary{listings,skins,breakable}
% 页眉页脚
\usepackage{fancyhdr}
\pagestyle{fancy}

% 定义草绿色
\definecolor{grassgreen}{RGB}{124,179,66}
% 补全缺失的浅蓝色定义（防止编译报错）
\definecolor{lightblue}{RGB}{30,144,255} 

% 清空默认
\fancyhf{}

% 页眉
\fancyhead[L]{\textcolor{lightblue}{\kaishu Environmental Ethics}}
\fancyhead[R]{\textcolor{lightblue}{\textit{2026.9-2027.1}}}

% 页脚页码居中
\fancyfoot[C]{\textcolor{lightblue}{\thepage}}

% 页眉、页脚线颜色
\renewcommand{\headrulewidth}{1pt}
\renewcommand{\footrulewidth}{1pt}
\renewcommand{\headrule}{\color{blue!30!white}\hrule height 1pt width\headwidth}
\renewcommand{\footrule}{\color{blue!30!white}\hrule height 1pt width\headwidth}
\newtcolorbox{fancybox}[2][blue]{ % 可选颜色参数
  enhanced,
  breakable,
    colback=#1!3!white,
    colframe=#1!18!white,
    boxrule=0.8pt,
    arc=6pt,
    outer arc=6pt,
  left=10pt,right=10pt,top=12pt,bottom=10pt,
  title={\textbf{#2}},
    fonttitle=\bfseries\color{#1!65!black},
  attach boxed title to top left={yshift=-2mm,xshift=5mm},
  boxed title style={
        colback=#1!8!white,
        colframe=#1!18!white,
        boxrule=0.6pt,
        arc=4pt,
  },
    shadow={1.5pt}{-1.5pt}{0pt}{black!10},
}

% 定义一套适合C++的配色方案
\usepackage{tabularx}
\newfontfamily{\consolasfont}{Ubuntu}

\newcommand{\bluecode}[1]{{\color{blue}\consolasfont #1}}
\lstdefinestyle{MyCppStyle}{
    language = C++,
    % 确保这里使用 \consolasfont
    basicstyle = \consolasfont\small, 
    keywordstyle = \color{blue!70!black}, 
    commentstyle = \color{green!50!black}, 
    stringstyle = \color{red!60!black}, 
    numberstyle = \tiny\color{gray}, 
    numbers = left, 
    frame = single, 
    rulecolor = \color{lightgray}, 
    breaklines = true, 
    backgroundcolor = \color{white}, 
    tabsize = 4, 
    showstringspaces = false, 
    xleftmargin = 2em, 
    escapeinside=``,
}
% --- lstdefinestyle HTML 风格定义 ---
\lstdefinestyle{mhtml}{
    language = HTML,
    basicstyle = \consolasfont\small, 
    tagstyle = \color{blue!70!black}\bfseries,
    keywordstyle = \color{purple!80!black},
    commentstyle = \color{green!50!black}\itshape, 
    stringstyle = \color{red!60!black}, 
    numberstyle = \tiny\color{gray}, 
    numbers = left, 
    frame = single, 
    rulecolor = \color{lightgray}, 
    breaklines = true, 
    backgroundcolor = \color{white}, 
    tabsize = 4, 
    showstringspaces = false, 
    xleftmargin = 2em, 
    escapeinside=``,
    morekeywords={als, doctype, html, head, title, meta, link, body, div, span, h1, h2, h3, h4, h5, h6, p, a, img, ul, ol, li, table, tr, th, td, pre, code, script, style},
    morekeywords=[2]{class, id, href, src, rel, type, charset, lang, name, content, style},
    keywordstyle=[2]\color{purple!80!black},
}
\usepackage{hyperref}
\hypersetup{
    colorlinks=true,      
    linkcolor=blue,       
    filecolor=magenta,    
    urlcolor=cyan,        
    pdftitle={我的文档标题}, 
}

% ==================== 正文部分 ====================
\begin{document}

% 开启无页眉页脚模式（针对封面和目录）
\pagestyle{empty}

% ------------------ 美化封面页 ------------------
\begin{titlepage}
    \centering
    % 顶部 TikZ 几何几何装饰条（全宽 bleed 效果）
    \begin{tikzpicture}[remember picture, overlay]
        \fill[blue!50!white] (current page.north west) rectangle ([yshift=-2cm]current page.north east);
        \fill[grassgreen] (current page.north west) ++(0,-2cm) rectangle ([yshift=-2.2cm]current page.north east);
    \end{tikzpicture}
    
    \vspace*{3.5cm}
    
    % 大标题
    {\fontsize{30pt}{36pt}\selectfont\bfseries\color{blue!80!black} 环境伦理学}\\[0.6em]
    % 副标题或英文标识
    {\fontsize{14pt}{18pt}\selectfont\color{gray}\textsf{Environmental Ethics}}\\[1cm]
    
    % 居中装饰线
    \centering\makebox[4cm]{\color{grassgreen}\rule{5cm}{1.5pt}}
    
    \vfill
    
    % 使用 tcolorbox 包装的精致作者信息框
    \begin{tcolorbox}[
        enhanced,
        width=0.55\textwidth,
        colback=blue!5!white,
        colframe=blue!5!white,
        boxrule=1.2pt,
        arc=6pt,
        sharp corners=downhill, % 别致的对角切角效果
        shadow={2pt}{-2pt}{0pt}{black!20},
        halign=center,
        valign=center,
        top=15pt, bottom=15pt
    ]
        {\large\bfseries\color{black!80} Author：} {\large\kaishu Simson} \\[0.6em]
        {\large\bfseries\color{black!80} Date：} {\large\kaishu September 2026}
    \end{tcolorbox}
    
    \vfill
    \vspace*{2cm}
    
    % 底部标识
    {\color{gray!80}\small\the\year\ \copyright\ All Rights Reserved.}
\end{titlepage}

\pagestyle{fancy}

\newpage
\tableofcontents
\newpage

% 设置从当前页（正文第一页）开始采用阿拉伯数字编码，并自动重置为第 1 页
\pagenumbering{arabic}

\section{概论}

环境伦理学研究环境相关的价值判断和伦理判断。相比于一般应用伦理学，要努力突破人类中心主义视角，对环境相关的伦理问题进行反思。环境伦理学的讨论分以下几个部分：

\begin{enumerate}
    \item 环境问题的历史和现状
    \begin{itemize}
        \item 环境问题的历史
        \item 20世纪中期依赖的环境危机
        \item 中国环境现状
    \end{itemize}
    \item 环境危机的根源
    \begin{itemize}
        \item 从经济，技术层面的反思
        \item 从文化价值层面的反思————对环境危机根源认识的深化
    \end{itemize}
    \item 环境伦理学的性质和起源
    \begin{itemize}
        \item 环境伦理学的学科性质
        \item 几种主要理论和观点
        \item 研究方法和研究思路
    \end{itemize}
    \item 伦理边界的扩展
    \begin{itemize}
        \item 伦理的进化
        \item 人类中心主义的主要形态和根据
        \item 非人类中心主义
        \item 两种立场
    \end{itemize}
    \item 自然的道德地位
    \begin{itemize}
        \item 利益原则与道德可考虑性（费因伯格）
        \item 树木具有法律地位吗（斯通）
    \end{itemize}
    \item 动物解放与动物权力论
    \begin{itemize}
        \item 动物的出境和人类道德
        \item Singer和动物解放：功利主义和道德眼神主义
        \item 动物权利论：义务论，天赋价值和生命主体
        \item 动物解放与素食主义
        \item 动物解放与环境伦理学的矛盾
    \end{itemize}
    \item 生物中心环境伦理学
    \begin{itemize}
        \item 生物中心论
        \item 罗尔斯顿
        \item 敬畏生命的伦理学（史怀哲）
        \item 尊重自然的环境伦理学（Paul Taylor）
    \end{itemize}
    \item 生态中心主义环境伦理学
    \begin{itemize}
        \item 生物中心主义环境伦理学的缺陷
        \item 利奥波德与大地伦理————整体主义伦理学
        \item 对整体主义环境伦理的批评
    \end{itemize}
    \item 深层生态学
    \begin{itemize}
        \item 奈斯：浅层的和深层的生态学
        \item 生态智慧与自我实现
        \item 生态哲学与宗教————精神生态学
    \end{itemize}
    \item 环境正义
    \begin{itemize}
        \item 起源：对非人类中心主义misanthropic倾向的批评
        \item 环境社会学的重要研究：环境污染与社会不公的关联
        \item 消除环境不公的根源：国内的和国际的
    \end{itemize}
    \item 社会生态学
    \begin{itemize}
        \item 马克思主义和生态社会主义
        \item 生态无政府主义————Murray Bucchin
        \item 生态乌托邦
    \end{itemize}
    \item 环境伦理的精神资源
    \begin{itemize}
        \item 中西方传统文化中支持环境伦理的人文，哲学和宗教思想
    \end{itemize}
    \item 走向实践的环境伦理
\end{enumerate}

\section{环境伦理学的发展与环境现状}

\subsection{环境问题}

环境问题的本质是人类与自然在物质和能量的交换上产生的问题。

\subsubsection{环境问题的历史}

\paragraph{环境史}
\begin{enumerate}
    \item 农耕时代：人的文化和社会是自然的循环的一部分。也有环境问题，但是和今天有本质的区别。
    \item 工业时代：利用化石燃料。城市的兴起令人类脱离自然
    \item 二战之后：化学合成物增多，人类或冻对地球的影响越来越大。
\end{enumerate}

在战后，用于武器的技术变化为核电站，DDT变化为农用杀虫剂。20世纪的环境危机体现在物流师年代开始各种公害事件频繁出现，比如：

\begin{itemize}
    \item LA光化学烟雾事件（50年代初）
    \item 日本水俣病（1956）
    \item 伦敦烟雾事件（1952）
\end{itemize}

七十年代以后，环境问题在范围上变得更广，危害人整体的生存和发展。水污染，大气污染，沙漠化和水土流失。关于环境治理的国际责任引发了复杂的环境正义问题的争端。

以上的历史现实催生了现代环境运动，和民权运动等相互结合，有以下几个重要的人物：

\paragraph{现当代环境运动}

\begin{itemize}
    \item 1962年雷切尔卡森《寂静的春天》，表达了对人通过技术来控制自然的批评。这是现代环境运动的先驱。
    
    思想谱系：继承史怀哲（序言）“敬畏生命”的观点，表达了对于人类中心主义的激烈批判。
    \item 1967历史学家霖林恩·怀特在《科学》发文“我们生态危机的历史根源”，指责基督教为生态危机负责。他揭示了基督教发展的历史和西方基督教文明对环境的破坏。
    \item 1968年经济学家提出“公地悲剧”，将环境问题拔高到经济学和社会学高度。（观点：需要靠胁迫性的共识，如法律、义务，形成纳税一类义务）
    
    \textbf{著名口号}：生育权不是一种人权，号召采取全球性政策阻止贫困国家大量生育。
    \item 1971年巴里·康芒纳《封闭的循环————自然，人和技术》，技术的生态缺陷来源于背后的科学方法。
    \item 1972年罗马俱乐部发表第一个研究报告《增长的极限》，语言经济增长不可能无限持续下去，因为石油等自然资源的供给有限，预测世界性灾难即将来临
    
    对策：提出“零增长”对策。（The limit to growth）
    \item 1973英国经济学家舒马赫在《小的是美好的》一书中对于西方经济学的环境后果做了分析和批评。
\end{itemize}

\begin{fancybox}{生态学四定律}
    \begin{enumerate}
        \item 每一种事物都与别的事物相关
        \item 一切事物都必然有其去向
        \item nature knows best
        \item 没有免费的午餐
    \end{enumerate}
\end{fancybox}

\subsection{环境与伦理学}

环境问题最终会变成一个道德问题。研究环境问题，可以有经济的观点，法律的观点，和道德的观点。环境哲学要对于各种观点进行比较。

伦理学是研究道德的学问。伦理学是对社会道德生活在理论上的概括和总结，他首先从理论上阐明认得本性，再阐明道德的来源，内容和总结，以及道德修养等问题。从而找到和指出关于真善美的路径。

进行道德判断的根据有两大流派：

\begin{itemize}
    \item 目的论或后果论：一件事情是否道德，要看其效果，产生的善是否超过恶（功利主义）。
    \item 非后果论（义务论）：有的行为本身就具有善的性质，人们做道德的事情是因为这件事情本身就是道德的。
\end{itemize}

规范伦理学的三个主题：价值论，义务论，德性论：

\begin{itemize}
    \item 价值论：什么样的事情或者行为是有价值的
    \item 义务论：我们对有价值的事物负有道德义务。
    \item 德性论：讨论我们的品德和人格理想。
\end{itemize}

道德的讨论不单纯的计算，还包含对于自然秩序的探讨。环境伦理不仅仅是应用伦理学，应用伦理学要求对相关的事实有清楚的了解，并进行合乎逻辑的道德推理。环境伦理学要吧人与人的道德关系推广和突破到是人与非人的生命甚至非生命存在物，就此而言环境伦理学是突破应用伦理学的界限的。

\subsection{传统环境伦理学如何处理环境问题}

把环境（动物，植物，土地，生态系统）看作财产，附属于人类，然后应用现有的伦理准则，人对环境只有间接义务。仁慈对待动物在欧洲有悠久的历史，出发点是人类的利益。但是存在缺陷：对于一些直觉上错误的行为无法有效进行谴责和制止。

当代的环境伦理也有CBA做法，试图以成本效益分析解决环境问题：

\begin{enumerate}
    \item 以货币等价，计算某项环境举措的成本和效益，使之最优化
    \item 如何评价环境的收益？健康和审美的价值应该如何计算？
    \item 推论：适度环保
\end{enumerate}

当然这种方法也有缺陷：经济目标优先，不平等（一美元一票而不是一人一票；发生了富人受益）；考虑不到对于后代的保护，比如核废水对于后代的经济负担（可能会有货币价格导致的通货膨胀，今天1美元，500年之后值160亿美元）

% 以下内容依据 Slides/2.西方环境伦理学（1）（伦理界限的扩大）.ppt 续写。
\section{西方环境伦理学：伦理边界的拓展}

\subsection{伦理观念的进化}

% 课件第2--5页，包括第3页的伦理观念进化图。
\subsubsection{道德共同体的扩大}

罗德里克·纳什（Roderick Nash）用伦理边界的逐步扩大来理解伦理观念的变化。受到道德关怀的对象，可以从自我、家庭、部落，扩大到地区、国家、种族和全人类；环境伦理学进一步追问，动物、植物、生命、岩石等无机物、生态系统乃至星球和宇宙，能否进入道德共同体。

这一图式提示我们：某类存在物在既有伦理中受到排斥，并不意味着这种排斥永远合理。不过，它首先是一种理解伦理扩展的思路，不能单凭扩展的历史趋势，就证明所有存在物都拥有相同的权利。每一次扩展仍然需要说明其道德根据。

课件也从法律史展示权利保护范围的变化，列举了《大宪章》、美国《独立宣言》、奴隶解放、妇女选举权、原住民公民资格、劳动者保护和黑人民权等例子，最后引向濒危物种保护。这里需要区分：\textbf{某类对象受到法律保护，与该对象自身成为权利主体，是两个不同的问题。}保护物种的法律，可以以人类的公共利益为根据，也可以引出自然是否拥有自身权利的进一步讨论。

\subsubsection{传统伦理中的关系范围}

课件以犹太教、基督教的摩西“十诫”和儒家的仁、忠恕之道为例，说明传统道德讨论主要围绕人与神、人与人之间的关系展开。“十诫”包含宗教义务以及尊敬父母、不可杀人、不可偷盗等规范；儒家强调“仁者爱人”和“己所不欲，勿施于人”。

环境伦理学由此提出：当行为直接影响动物、植物和生态系统时，我们是否只需考虑受到影响的人？如果自然也具有道德地位，那么道德义务的对象就需要进一步扩大。课件的这些例子用于说明传统伦理的讨论重心，不能据此穷尽这些传统中有关自然的全部思想。

\subsection{人类中心主义及其不同形态}

% 课件第6--8页。
\subsubsection{人类中心主义的几种含义}

人类中心主义（anthropocentrism）在不同层面具有不同含义，讨论时需要先明确所指：

\begin{enumerate}
    \item \textbf{宇宙论层面}：人类如何理解自身在宇宙中的位置。课件以地心说、日心说所涉及的世界图景为讨论背景。
    \item \textbf{生物学层面}：把人看作最高的存在物，并认为自然界的其他存在物为人类而存在。
    \item \textbf{哲学和伦理学层面}：把人看作自然的主宰，以人的价值和利益作为评价自然、决定如何对待自然的最终尺度。
\end{enumerate}

课件列举的思想线索包括古希腊“人是万物的尺度”、基督教关于人具有“神的形象”的观念，以及近代哲学中培根对知识和力量的强调、笛卡尔对思维主体的强调、康德关于人是目的的思想。这些线索共同引出人的特殊地位问题，但不同思想家的论证及其对自然的态度仍有差别。

\subsubsection{西方传统中对自然和动物的态度}

\begin{itemize}
    \item \textbf{基督教传统中的不同方向}：一方面存在征服、支配自然的理解，另一方面也有通过自然赞美上帝的理解。圣方济各（St. Francis of Assisi）强调人与动物同在上帝面前，体现了另一种对待自然的精神资源。
    \item \textbf{托马斯·阿奎那}：以理智和理性说明人的自由意志及责任能力，并以人与动物在这些能力上的差异，为动物供人使用提供根据。
    \item \textbf{笛卡尔}：课件以“动物是机器”概括其动物观。这一观点削弱了把动物作为有自身感受和利益的对象来考虑的理由。
    \item \textbf{康德}：课件强调人对动物只有间接义务。善待动物的要求通过人类自身的道德义务得到说明，动物尚未成为直接义务的对象。
\end{itemize}

\subsubsection{强人类中心主义与弱人类中心主义}

\paragraph{强人类中心主义}
课件将其概括为认可人的一切偏好，把满足既有欲望作为行动根据。在环境问题上，它容易表现为经济中心主义，把人的幸福主要换算为经济产值，以经济增长衡量自然利用的成效。

\paragraph{弱人类中心主义}
弱人类中心主义强调全人类的利益以及人类的长远利益，要求超出特定个人、群体和当代人的眼前偏好。它也承认自然具有审美、精神等价值，因而保护自然的理由可以超出直接的物质收益。

两者的分歧在于如何理解人的利益，以及是否应当反思和约束人的偏好。弱人类中心主义可以支持更严格、更长远的环境保护，但其最终的价值根据仍然是人类的利益。

\subsection{对人类中心主义的评价与物种利他主义}

% 课件第9--10页。
\subsubsection{认识立场、物种身份与价值判断}

课件参照杨通进的观点，将人类中心主义分为三个层次：

\begin{enumerate}
    \item \textbf{认识论层次}：人类的认识具有人的视角，我们通过自身的感知、概念和语言理解世界。课件把这一层次视为无法摆脱的认识条件。
    \item \textbf{生物物种层次}：人类属于生物界，受到生存条件和生态规律的约束。课件同样把这一点视为无法取消的事实。
    \item \textbf{价值论层次}：认为只有人具有内在价值，其他存在物的价值完全取决于它们对人的用途。这是一种可以受到批评、也可能被超越的价值立场。
\end{enumerate}

因此，\textbf{人只能从人的认识条件出发，并不能直接推出只有人才值得道德关怀。}环境伦理学对人类中心主义的批评，主要针对价值论上的排他性。人在生物学上属于某个物种，也不能单独构成只顾本物种利益的道德理由。

\subsubsection{利他主义的争论}

课件把人类中心主义与个人利己主义作类比：正如人可能只承认自身利益，也可能只承认本物种的利益。要讨论能否超出物种界限，可以先考察个人能否超出自身利益。

关于利他行为，课件列出几种解释：

\begin{itemize}
    \item 人能够出于同情、移情或宗教信念，为他人牺牲自己的利益。
    \item 从“自私的基因”等思路出发，可以尝试为表面的利他行为寻找生物学解释。
    \item 心理利己主义认为，帮助别人最终仍带来精神满足，因此行为最终仍以自身为目的。
\end{itemize}

这里要区分行为的后果与行为的动机：帮助他人后感到满足，并不足以单独证明，获得满足就是帮助他人的唯一目的。同样，对行为作生物学解释，也不能直接决定这种行为在道德上应当如何评价。

\subsubsection{物种利他主义的可能性}

物种利他主义（species altruism）追问：人能否为了其他物种自身的利益，约束甚至牺牲本物种的部分利益？如果承认这一可能性，保护动物和自然就可以出于对其自身处境的关切。

课件引用马克思《1844年经济学哲学手稿》关于生产尺度的论述：动物按照所属物种的尺度和需要来建造，人则能够理解不同物种的尺度，并把对象自身的尺度运用到对象上。这里的启发是，人具有理解其他存在物的生存方式、反思自身需求的能力。这为跨越物种界限的道德关怀提供了思路，具体义务仍需进一步论证。

\subsection{自然的道德地位：费因伯格的利益原则}

% 课件第11--31页；英文原文与中文译文合并整理。
\subsubsection{从环境人权到自然自身的权利}

课件以1971年在佐治亚大学召开的环境哲学会议，以及1974年出版的《哲学与环境危机》为早期环境哲学讨论的例子。威廉·布莱克斯通（William Blackstone）提出，适于生存的环境是一项基本人权。环境条件的改变要求人们根据公共福利和公平，限制过去被视为当然的自由，特别是财产权的行使。

这一立场能够说明环境保护为何可以限制个人利用资源和排放污染的自由，但论证仍以人的权利为中心。约尔·费因伯格（Joel Feinberg）进一步追问：\textbf{哪些存在物能够成为权利的拥有者？}他的回答以“利益原则”为核心。

\begin{fancybox}{理解自然道德地位的几个概念}
    \begin{itemize}
        \item \textbf{道德可考虑性}：在判断行为时，某个存在物自身的处境是否应当得到考虑。
        \item \textbf{道德权利}：是否存在可以为该存在物自身提出的正当要求，并由此对他人形成义务。
        \item \textbf{法律地位}：法律制度是否允许以该对象的名义提出请求，并对其所受损害给予救济。
    \end{itemize}
    这几个问题相关，但不能仅凭其中一个问题的答案，就直接确定其余问题的答案。
\end{fancybox}

\subsubsection{对他人有价值与拥有自身的利益}

费因伯格以泰姬陵等珍贵对象为例，区分两种保护理由。人们可以因为文化遗产、民族情感或审美欣赏而承担维护它的义务；如果它遭到破坏，损失可能非常重大。然而，这些理由首先说明它对人有价值，尚未说明它自己具有能够得到满足或受到挫折的利益。

按照课件介绍的这一立场，“纯粹的物”缺少欲望、希望、冲动和自身的目的，因而缺少属于它自己的福利或善。对物的保管义务，可以是对所有者、公众或后代承担的义务，并不必然对应物自身的权利。这里所说的物自身的权利，也不应与法律上人对物享有的“物权”混同。

\subsubsection{利益原则的论证}

费因伯格认为，能够拥有权利的存在物必须具有自身的利益，其理由有两个方面：

\begin{enumerate}
    \item \textbf{能够被代表}：他人能够代表权利主体提出要求，前提是存在属于该主体的利益。没有这种利益，就无法说明代理人究竟在为谁、为了什么提出要求。
    \item \textbf{能够成为受益者}：权利保护必须能够使该主体自身受益，或者避免它自身受到伤害。如果一个对象没有属于自己的善，就难以说它是权利保障的受益者。
\end{enumerate}

因此，拥有自身利益是拥有权利的必要条件。这个结论本身还没有证明，只要拥有某种利益，就自动拥有一切相应的权利；具体权利的内容和限度仍需讨论。

\subsubsection{动物为什么可能拥有权利}

课件以照料受遗赠资金保护的猫狗为例。照料者承担的职责可以直接以动物的利益为对象：动物有食欲、冲动和基本目的，其需要得到满足或受到挫折，会影响它自身的福利。因此，动物可以成为照料行为的受益者。

这一点也说明，\textbf{能够承担道德责任，与能够受到道德保护，所要求的能力不同。}动物即使不能成为对自身行为承担道德责任的行为主体（moral agent），仍可能有值得保护的利益。不能自行表达权利要求，也不必然取消被代表和受到保护的可能性。

费因伯格的论证首先说明，对至少某些动物谈论权利具有意义。它并未直接决定每一种动物都有什么具体权利，也没有取消讨论利益冲突的必要。

\subsubsection{植物的生长需要能否构成利益}

费因伯格承认树木具有生长和发展的自然趋向，却认为这种趋向不足以构成权利所需的利益。按照课件呈现的论证，利益至少以基本的认知能力为前提；树木缺少有意识的需求和目标，也没有满足、挫折、快乐或痛苦的体验。

围绕植物的争论，需要区分“需要”的两种用法：

\begin{enumerate}
    \item \textbf{实现功能的条件}：例如，植物需要水分，意思是水分是其生长或完成某种生理功能的必要条件。这首先是关于功能和因果关系的事实判断。
    \item \textbf{对主体自身有好处}：说某物符合一个主体的利益，意味着它有助于该主体自身的福利，缺少它会使该主体受到伤害。这包含关于什么对该主体是好或坏的价值判断。
\end{enumerate}

费因伯格拒绝仅凭第一种需要，就直接推得第二种意义的利益。日常语言中，我们也会说植物“茁壮”“繁荣”或“萎靡”。这些说法能否证明植物拥有自身的善，正是争论所在，不能只凭用词就完成论证。

由此留下的问题是：\textbf{利益是否必须依赖认知和主观体验？}如果认为一个生命的生长、完整和繁衍本身就构成对它的好处，就需要为这种更宽的利益概念提供根据。这也是进一步讨论生物中心主义的重要问题。

\subsubsection{物种与其他边界案例}

课件介绍，费因伯格也把物种本身排除在权利主体之外。物种与具有感受、欲望的个体不同；公司、国家之所以能够被代表，则与其组织结构和对成员利益的代表有关，不能仅凭它们具有法律人格，就直接推得物种具有相同的地位。

课件还把死者、持续无意识状态者、婴儿和后代列为检验利益原则的边界案例，并介绍了其中排除死者及所谓“植物人”权利的主张。这些是课件中需要评估的理论结论。就论证而言，需要进一步说明：过去的愿望、未来可能形成的利益，以及意识状态的差异，应当如何纳入利益原则。不能只用“当下能否表达利益”替代这些问题。

利益原则由此扩大了对动物的道德考虑，同时也形成了新的边界。它能否充分说明植物、物种和生态整体的价值，是环境伦理学继续推进的方向。

\subsection{树木具有法律地位吗：斯通的主张}

% 课件第32--36页。
\subsubsection{矿王谷开发争议与诉讼资格}

课件以塞拉俱乐部诉莫顿案（Sierra Club v. Morton）引入自然物的法律地位问题。争议围绕迪斯尼拟在矿王谷（Mineral King）进行的开发及其对当地审美和生态环境的影响展开。案件的重要问题是，环保组织是否具有把这一争议提交法院的诉讼资格（standing）。

这里需要区分：某项开发是否损害环境，以及由谁、根据何种自身受损的理由请求法院审查。课件记载，1972年最高法院判塞拉俱乐部败诉，其关键障碍在于诉讼资格。环境可能受到损害，并不足以自动解决谁可以起诉的问题。

\subsubsection{自然物作为权利主体的三个条件}

克里斯托弗·斯通（Christopher D. Stone）在《树木具有法律地位吗？——论自然物的法律权利》中，提出赋予森林、海洋、河流等自然物法律权利。其主张并不要求自然物拥有与人完全相同的权利体系，也不要求所有自然物都拥有完全相同的权利。

按照课件的概括，自然物真正具备法律地位，需要满足三个条件：

\begin{enumerate}
    \item \textbf{能够以自身名义寻求司法保护}：请求可以为了自然物自身而提出。自然物无法亲自表达或起诉的问题，可以通过监护人、代表人制度处理。
    \item \textbf{自身所受伤害得到考虑}：法院确定救济时，需要考虑自然物受到的损害，而不只计算人类财产或使用利益的损失。
    \item \textbf{救济为了其自身利益而实施}：救济应当用于保护、恢复相关自然物，使其成为救济的受益对象。
\end{enumerate}

斯通认为，“河流享有法律权利”这样的表述能够影响制度的发展。它会改变谁能够提出请求、哪些损害进入司法考虑，以及救济如何安排。因此，赋予自然物法律权利具有制度上的实际含义。

\subsubsection{自然的代表与利用的限度}

斯通并不由赋予自然权利推出全面停止人类利用自然。承认权利，需要同时讨论利用的条件、边界以及受到损害后的救济。其目标是使自然物在法律制度中获得更有力的地位，让有关其自身处境的要求得到表达。

课件介绍，道格拉斯法官支持自然物可以有法律代言人的思路，认为与受到威胁的自然对象关系密切的人可以承担代表职责。理解这一点时，应当将这一意见与案件多数意见区分。课件另记载，迪斯尼于1978年放弃开发计划；这一后续结果也不能等同于最高法院已经承认自然物的权利。

\subsubsection{尚待解决的问题}

\begin{itemize}
    \item \textbf{如何确定自然的利益？}如果不同的管理方式分别有利于某些个体、物种或生态过程，需要说明以什么作为评价标准。
    \item \textbf{谁能够代表自然？}荒野保护组织、木材公司和开发者可能都声称了解自然，但它们具有不同的利益和目标。与自然关系密切可以成为代表资格的考虑因素，仍需进一步判断其代表是否恰当。
    \item \textbf{道德根据与法律安排如何衔接？}费因伯格主要追问何种存在物具有自身利益、因而可能拥有权利；斯通主要追问如何通过代表、诉讼与救济制度，使自然得到以自身为对象的法律保护。两种讨论相互关联，同时保留各自需要回答的问题。
\end{itemize}

\section{土地伦理}

土地伦理规范：不管是不是有经济利益可图，就生物权利而言，土地这个生态共同体的成员应该存在。

\begin{fancybox}{Decision}
    当一件事情倾向于保证生物共同体的完整，稳定和美感时，它就是正当的，反之是不正当的。
\end{fancybox}

生态良心：如果没有良知，义务就没有任何意义，而我们面临的问题就是将社会良知从人扩展到土地。

\subsubsection{大地伦理学的当代阐释}

\begin{itemize}
    \item 伦理学的基础：道德情感，同情共情，道德情操论
    \item 生态学的形而上学意蕴：整体论哲学
    \item 个体与整体：痛苦有进化价值，不是“恶”
\end{itemize}

西方的整体论传统：柏拉图，黑格尔。

在西方古典的更为宽广的伦理哲学背景中，大地伦理是完全可以理解的。所以关于柏拉图关于公共和私人正义的体系是一种恰当的伦理体系

\subsection{对整体主义环境伦理学的批评和辩护}

两个古典的问题：

\begin{itemize}
    \item 整体有道德资格吗？
    \item 如何从生态学的是到伦理学的应该？
\end{itemize}

生态学的描述蕴含价值？生态学模型：有机体organisms-群落community-能量energy。

还有更加尖锐的批评：Tom Regan指责（生态法西斯主义）

为了更大的生态的善而牺牲个体，很难看出个体权利在其中所处的位置。

Zimmerman：政府可能以保护环境的名义没收私有财产，对人的自由造成压迫。

\subsubsection{辩护}

同心圆模型：整体主义伦理学是对原来伦理学的一种增量而不是代替。

二阶原则：

\begin{enumerate}
    \item 依据远近亲疏来决定优先性
    \item 依据利益的强弱决定优先性
\end{enumerate}

赫夫南的修正：一个行为，当它有助于维护生态系统的完整，稳定和美丽时，它就是正确的，假如我这样做时候不犯其他错误的话。【加入限定条件：不违背重要的人际义务】

\subsection{生态学}

\subsubsection{生态学的发展历史}

神学生态学：英国植物学家John Ray，《呈现于创造作品中的上帝的智慧》描述了物种丰富，自然和谐，体现了上帝的特性。

林奈“以自然的安排称呼自然的和谐有序”，虽然他对自然的认识是机械性和局限性的。

1866年，海克尔首次发明生态学的英语词汇

美国生态学家Clements，在20世纪明确提出要研究生态的系统的社群。将生态学学科化，科学化。

最后形成了生态学的学科体系和研究方法的概念。

\subsection{生态哲学基础的本体论}




\end{document}
